\subsection{The universal Schur coefficient map}

The distinction between an isotypic component and a specified line in
its multiplicity space is essential here.  We establish the coefficient
identity before fixing the web.  All vector spaces in this subsection
are complex and four-dimensional.  Put
\[
 \mathcal F(X)=\bigwedge^6\Sym^2X,\qquad
 \Lambda=\{(5,4,2,1),(4,4,4,0)\}.
\]
For a linear map \(T\) write \(\mathcal F(T)=\bigwedge^6\Sym^2T\).

\begin{lemma}[Universal coefficient identity]
\label{lem:universal-schur-readout}
The natural degree-twelve map
\begin{equation}\label{eq:universal-coefficient-map}
 \begin{split}
 \Phi_{V,U}:\mathcal F(V)^*\otimes\mathcal F(U)
       &\longrightarrow\Sym^{12}(V^*\otimes U),\\
 \Phi_{V,U}(\beta\otimes\xi)(T)&=\beta(\mathcal F(T)\xi)
 \end{split}
\end{equation}
is \(GL(V)\times GL(U)\)-equivariant.  Under the functorial
multiplicity-free decompositions
\begin{equation}\label{eq:minor-two-summands}
 \mathcal F(X)=\mathbb S_{(5,4,2,1)}X
                    \oplus\mathbb S_{(4,4,4,0)}X
\end{equation}
and the Cauchy decomposition
\[
 \Sym^{12}(V^*\otimes U)
   =\bigoplus_{\substack{\lambda\vdash12\\\ell(\lambda)\le4}}
       \mathbb S_\lambda V^*\otimes\mathbb S_\lambda U,
\]
the restriction of \(\Phi_{V,U}\) to
\(\mathbb S_\lambda V^*\otimes\mathbb S_\mu U\) is zero
if \(\lambda\ne\mu\).  For \(\lambda=\mu\in\Lambda\),
it is a nonzero scalar multiple of the canonical Cauchy inclusion.
With compatible functorial normalizations this scalar is one.
Consequently, for every fixed \(\beta=\sum_\lambda\beta_\lambda\),
\begin{equation}\label{eq:fixed-slice-universal}
 \Phi_{V,U}(\C\beta\otimes\mathcal F(U))
    =\bigoplus_{\lambda\in\Lambda}
          \C\beta_\lambda\otimes\mathbb S_\lambda U,
\end{equation}
where each tensor product on the right is placed in its Cauchy summand.
A zero component contributes the zero subspace.
\end{lemma}
\begin{proof}
The action on matrices is \(T\mapsto gTh^{-1}\); the action on
functions is \(f(T)\mapsto f(g^{-1}Th)\).  Functoriality of
\(\mathcal F\) shows directly that \eqref{eq:universal-coefficient-map}
intertwines this action with
\(\beta\otimes\xi\mapsto(\mathcal F(g^{-1})^*\beta)
\otimes\mathcal F(h)\xi\).  Thus equivariance is a statement about
the universal map, not about a fixed-\(\beta\) slice.

For completeness, the exterior character identity used here is
\[
 [z^6]\prod_{1\le i\le j\le4}(1+zx_ix_j)
       =s_{(5,4,2,1)}(x)+s_{(4,4,4,0)}(x).
\]
Equivalently, it follows by exterior duality from
\eqref{eq:two-web-components} and
\(\det(\Sym^2X)=(\det X)^5\).  The two dimensions are
\(175\) and \(35\), whose sum is \(\binom{10}{6}\).
The characteristic-zero Schur decomposition is functorial; the
Cauchy formula is multiplicity-free as a representation of the
product group \cite{FultonHarris,DCEP}.  The source of \(\Phi\)
is the direct sum of the four irreducible product representations
indexed by \(\Lambda\times\Lambda\).  Only the two diagonal
ones occur in the target.  Schur's lemma kills the off-diagonal
blocks and makes each diagonal block a scalar multiple of its
Cauchy inclusion.  This also proves that no other left multiplicity
line can enter the formula.

Here is an explicit nonvanishing check for both scalars.  For an
ordered basis \(x_1,\ldots,x_4\), order the monomials of
\(\Sym^2X\) as
\[
 (b_1,\ldots,b_{10})=
 (x_1^2,x_1x_2,x_1x_3,x_1x_4,x_2^2,x_2x_3,
                  x_2x_4,x_3^2,x_3x_4,x_4^2).
\]
Highest-weight vectors for the two summands are
\[
 \xi_{5421}=b_1\wedge b_2\wedge b_3\wedge b_4\wedge b_5\wedge b_6,
 \qquad
 \xi_{4440}=b_1\wedge b_2\wedge b_3\wedge b_5\wedge b_6\wedge b_8.
\]
Their weights are respectively \((5,4,2,1)\) and \((4,4,4,0)\).
Every raising operator \(x_i\partial/\partial x_j\), \(i<j\),
either gives zero or repeats a wedge factor, so both are highest
vectors.  Complete reducibility and the dimension sum above show that
these two generated irreducible submodules exhaust \(\mathcal F(X)\);
this also proves \eqref{eq:minor-two-summands} directly.
Choose \(\beta_\lambda\) in the dual summand with
\(\beta_\lambda(\xi_\lambda)=1\), and extend it by zero on the
other summand.  After identifying the ordered bases of \(U\) and
\(V\), evaluation at \(T=\diag(t_1,t_2,t_3,t_4)\) gives
\begin{equation}\label{eq:two-block-nonvanishing}
 \begin{split}
 \Phi(\beta_{5421}\otimes\xi_{5421})(T)&=t_1^5t_2^4t_3^2t_4,\\
 \Phi(\beta_{4440}\otimes\xi_{4440})(T)&=t_1^4t_2^4t_3^4.
 \end{split}
\end{equation}
In particular both evaluations are one at the identity.  This proves
nonvanishing without a computed rank or a choice of numerical web.
The canonical coefficient realization of the Cauchy inclusion is
\(\beta\otimes\xi\mapsto[T\mapsto\beta(\mathbb S_\lambda(T)\xi)]\).
Compatible functorial splittings identify the restriction of
\(\mathcal F(T)\) with \(\mathbb S_\lambda(T)\), giving scalar one;
other scalar normalizations have the same image.  Finally let \(\xi\)
vary independently in its two summands.  The off-diagonal vanishing
and the nonzero diagonal scalars give \eqref{eq:fixed-slice-universal}.
\end{proof}

\begin{lemma}[Exterior duality and the common determinant twist]
\label{lem:readout-exterior-twist}
Set \(W=\Sym^2V\).  The canonical isomorphism
\begin{equation}\label{eq:explicit-exterior-duality}
 \begin{split}
 \bigwedge^4W\otimes(\det W)^*&\longrightarrow\bigwedge^6W^*,\\
 w\otimes\omega&\longmapsto[\eta\mapsto\omega(w\wedge\eta)]
 \end{split}
\end{equation}
sends the Pluecker line of \(R=\ker\gamma_R\) to the line of
\(\beta_R=\gamma_R^*(\operatorname{vol}_{S_R})\).
It identifies \(E_{175}\otimes(\det V)^{-5}\) with
\(\mathbb S_{(5,4,2,1)}V^*\), and
\(E_{35}\otimes(\det V)^{-5}\) with
\(\mathbb S_{(4,4,4,0)}V^*\).
\end{lemma}
\begin{proof}
Extend a basis of \(R\) to a basis of \(W\).  Both sides of the
first assertion are nonzero alternating six-forms annihilating a
wedge containing a vector of \(R\), and on the six complementary
basis vectors they agree after one scalar normalization.  This proves
the line identity.  The map in \eqref{eq:explicit-exterior-duality}
is natural for all isomorphisms of \(W\).
The highest weights of the two twisted summands are
\((-1,-2,-4,-5)\) and \((0,-4,-4,-4)\), respectively: these are
the negatives of the reversed partitions \((5,4,2,1)\) and
\((4,4,4,0)\).  Thus both labels and the common twist are as stated.
Changing the volume multiplies both components by the same scalar.
\end{proof}
